%%% FIELD proof-star-regression
Put \(a=\mathbf 1_{A_z}\).  By \(\texttt{def:binary-star-library}\),
\[
 \Sigma^{\star}_{qq}=1\ (q\ne z),\qquad
 \Sigma^{\star}_{zz}=1+|A_z|,\qquad
 \Sigma^{\star}_{zq}=\Sigma^{\star}_{qz}=a_q,
\]
and all other off-diagonal entries vanish.  If \(j=z\), then \(P\subseteq V\setminus\{z\}\), so
\[
 \Sigma^{\star}_{P,P}=I_{|P|},\qquad \Sigma^{\star}_{P,z}=a_P.
\]
The definition of \(r_j\) gives \(r_q=\mathbf 1\{q\in A_z\}\).

Suppose \(j\ne z\).  If \(z\notin P\), then \(\Sigma^{\star}_{P,j}=0\), because \(j\notin P\) and distinct noncentral coordinates are uncorrelated.  If \(j\notin A_z\), the same conclusion holds even when \(z\in P\), because \(\Sigma^{\star}_{zj}=0\).  Positive definiteness of \(\Sigma^{\star}_{P,P}\) therefore gives \(r=0\) in both cases.

It remains that \(j\in A_z\) and \(z\in P\).  Set \(Q=P\setminus\{z\}\), \(c=\mathbf 1_{A_z\cap Q}\), and order \(P\) as \(Q,z\).  Then
\[
 \Sigma^{\star}_{P,P}=
 \begin{pmatrix}I_{|Q|}&c\\ c^{\mathsf T}&1+|A_z|\end{pmatrix},
 \qquad
 \Sigma^{\star}_{P,j}=\binom{0}{1}.
\]
For
\[
 d=1+|A_z|-|A_z\cap Q|=1+|A_z\setminus Q|,
\]
direct multiplication gives
\[
 \begin{pmatrix}I_{|Q|}&c\\ c^{\mathsf T}&1+|A_z|\end{pmatrix}
 \binom{-c/d}{1/d}=\binom{0}{1}.
\]
Thus \(r_z=1/d\), \(r_q=-1/d\) on \(A_z\cap Q\), and \(r_q=0\) on \(Q\setminus A_z\).  The Schur-complement determinant identity gives
\[
 \det(\Sigma^{\star}_{P,P})=1+|A_z|-c^{\mathsf T}c=d.
\]
Here \(j\in A_z\setminus Q\), so \(2\le d\le1+|A_z|\le p\); in particular, every division above is valid.

%%% FIELD proof-separation
We first prove that the certificate formats vanish on the rejected model.  Fix a rejected state \(J\) and a topological order \(\pi_J\).  If \(u,v\) are nonadjacent, \(u\) precedes \(v\), and \(S=\operatorname{Pred}_{\pi_J}(v)\setminus\{u\}\), then \(\operatorname{pa}_{G_J}(v)\subseteq S\).  In the acyclic SEM of \(\texttt{def:covariance-model}\), the innovation of \(X_v\) is independent of all predecessor variables and its structural mean is measurable with respect to \(X_S\).  Hence
\[
 \operatorname{Cov}(X_u,X_v\mid X_S)=0.
\]
Positive definiteness and the Schur-complement formula yield
\[
 f_{uvS}(\Sigma)=\det(\Sigma_{S,S})
 \bigl(\Sigma_{uv}-\Sigma_{uS}\Sigma_{S,S}^{-1}\Sigma_{Sv}\bigr)=0,
\]
where \(\det(\Sigma_{\varnothing,\varnothing})=1\).  Likewise, for \(P=\operatorname{pa}_{G_J}(j)\), the innovation of \(X_j\) is independent of \(X_P\), so \(r_j(P;\Sigma)\) is its structural parent-coefficient vector.  The block constraint gives \(r_u=r_v\) for \(u,v\) in one block, whence
\[
 L_{j;u,v}(\Sigma)=\det(\Sigma_{P,P})(r_u-r_v)=0.
\]
These arguments use acyclicity, positive innovation variances, and the parent-block restriction, with no faithfulness assumption.

Assume now that \(\mathsf S(C)\ne\mathsf S(D)\).  The cases are disjoint because the router uses the first mismatch.

\emph{Skeleton mismatch.}
Choose an adjacency \(\{a,z\}\) present only in a source state \(E\), oriented \(a\to z\), and activate the block of \(\Pi^E_z\) containing \(a\).  Let \(J\) be the other state.  Order the nonadjacent endpoints as \(u,v\) in a topological order of \(J\), and put \(S=\operatorname{Pred}_{\pi_J}(v)\setminus\{u\}\).  Neither endpoint is in \(S\), and \(\Sigma^{\star}_{S,S}=I\).  Distinct active leaves are independent, so conditioning on active leaves in \(S\) does not change \(\operatorname{Cov}(X_a,X_z)=1\).  Therefore
\[
 f_{uvS}(\Sigma^{\star})=1.
\]

\emph{Collider mismatch.}
With equal skeletons, choose \(a\to z\leftarrow c\) as an unshielded collider only in \(E\), and activate the one or two blocks at \(z\) containing \(a,c\).  In any topological order of the noncollider state \(J\), the center \(z\) precedes the later endpoint \(v\) of \(a,c\), in each possible fork or chain orientation.  Let \(u\) be the earlier endpoint, \(S=\operatorname{Pred}_{\pi_J}(v)\setminus\{u\}\), and \(R=S\setminus\{z\}\).  Then \(a,c\notin S\), \(z\in S\), and projection of \(X_z\) on \(X_R\), using \(\texttt{lem:binary-star-regression-calculus}\), gives
\[
 d=1+|A_z\setminus R|,\qquad
 \det(\Sigma^{\star}_{S,S})=d,\qquad
 \operatorname{Cov}(X_a,X_c\mid X_S)=-d^{-1}.
\]
Thus \(f_{uvS}(\Sigma^{\star})=-1\).

\emph{Nonsingleton-family mismatch.}
For \(B\in\Pi^E_j\), define the edge atom
\[
 \mathscr E_j(B)=\{\{i,j\}:i\in B\}.
\]
These atoms partition the skeleton edge set.  An atom with at least two edges has a unique common center, so it recovers its target and nonsingleton block.  Thus equal skeletons but unequal vertex-indexed nonsingleton families give unequal edge-atom partitions.  For two unequal finite partitions, after possibly interchanging their roles, there are a source atom \(U\) and rejected atom \(K\) with
\[
 U\cap K\ne\varnothing,\qquad K\setminus U\ne\varnothing;
\]
otherwise the two partitions refine one another.  Activate the single source block represented by \(U\), with center \(z\), and let \(K\) have center \(j\).

If \(j=z\), choose rejected-block parents \(u,v\) represented by edges in \(U\cap K\) and \(K\setminus U\).  Then the regression lemma gives \(r_u=1\), \(r_v=0\), and \(\det(\Sigma^{\star}_{P,P})=1\), so \(L_{j;u,v}(\Sigma^{\star})=1\).  If \(j\ne z\), the two star atoms meet only in \(\{j,z\}\).  Hence \(j\to z\) in \(E\), \(z\to j\) in \(J\), and \(z\) is in \(K\).  For any other parent \(v\in K\), the regression lemma gives
\[
 r_z=d^{-1},\qquad r_v\in\{0,-d^{-1}\},\qquad
 \det(\Sigma^{\star}_{P,P})=d,
\]
so \(L_{j;z,v}(\Sigma^{\star})\in\{1,2\}\).

\emph{Family-mismatch orientation.}
The preceding invariants now agree, but \(F(C)\ne F(D)\).  There is an adjacency oriented \(x\to y\) in one state and \(y\to x\) in the other with \(\mathcal N_x\ne\mathcal N_y\).  This edge is a singleton block at both receiving endpoints; a nonsingleton occurrence would contradict equality of the corresponding vertex-indexed family.

Take the state with \(x\to y\).  If \(S\in\mathcal N_x\) and \(s\in S\), then \(s\to x\) occurs in both states.  If \(s,y\) were nonadjacent, reversing \(x-y\) would change \(s\to x\to y\) into the unshielded collider \(s\to x\leftarrow y\), contrary to equality of collider sets.  They are adjacent, and acyclicity rules out \(y\to s\) because \(s\to x\to y\to s\) would be a cycle.  Hence
\[
 S\subseteq\operatorname{pa}_{G_E}(y).
\]
Interchanging states and endpoints gives \(T\subseteq\operatorname{pa}_{G_J}(x)\) for every \(T\in\mathcal N_y\).

Some tail block is split by the source partition at the head in at least one direction.  Otherwise every \(S\in\mathcal N_x\) lies in a \(T\in\mathcal N_y\), and every such \(T\) lies in a \(U\in\mathcal N_x\).  Then \(S\subseteq T\subseteq U\); since \(S,U\) intersect and are blocks of the same partition, \(S=U\), so \(S=T=U\).  Applying the symmetric argument to every \(T\) gives \(\mathcal N_x=\mathcal N_y\), a contradiction.

Rename the selected direction so that \(x\to y\) is in \(E\), \(S\in\mathcal N_x\) is split by \(\Pi^E_y\), and \(T\in\Pi^E_y\) meets but does not contain \(S\).  Activate the distinct blocks \(\{x\}\) and \(T\) at \(y\).  In \(J\), regress \(x\) on \(P=\operatorname{pa}_{G_J}(x)\), set \(Q=P\setminus\{y\}\), and choose \(s\in S\cap T\), \(t\in S\setminus T\).  The common block \(S\) is contained in \(Q\).  Since \(A_y=\{x\}\cup T\), the regression lemma yields
\[
 d=1+|A_y\setminus Q|=2+|T\setminus Q|,\qquad
 r_s=-d^{-1},\qquad r_t=0,\qquad
 \det(\Sigma^{\star}_{P,P})=d.
\]
Thus \(L_{x;s,t}(\Sigma^{\star})=-1\).  The term \(T\setminus Q\) explicitly covers extra parents, and other reversed arrows have zero coefficient in this star.

The four cases exhaust unequal signatures and activate at most two whole blocks at one target.  Set all other source coefficients to zero and all innovation variances to one.  Since \(z\notin A_z\),
\[
 (B^{\star})^2=e_z(\mathbf 1_{A_z}^{\mathsf T}e_z)\mathbf 1_{A_z}^{\mathsf T}=0,
 \qquad (I_p-B^{\star})^{-1}=I_p+B^{\star}.
\]
Thus the constructed covariance is the displayed \(\Sigma^{\star}\in\mathcal M(E)\).  Its nonzero certificate and the universal vanishing proved above imply \(\Sigma^{\star}\notin\mathcal M(J)\).  In a topological order \(I_p+B^{\star}\) is unit triangular, so
\[
 \det(\Sigma^{\star})=\det(I_p+B^{\star})^2=1.
\]
Its entries are \(0,1\), and the center variance \(1+|A_z|\le p\).  Every recovered coefficient is \(0\), \(1\), or \(\pm1/d\), with \(2\le d\le p\) whenever division occurs; the cleared certificate values above lie in \(\{-1,1,2\}\).  Hence every regression numerator is in \(\{-1,0,1\}\), every denominator and every other numerical magnitude is at most \(p\), and the height is at most \(B(p)=\lceil\log_2p\rceil\).  For \(p=1\) the implication is vacuous.

Finally, skeleton comparison, collider enumeration, atom comparison, containment scans, topological sorting, and elimination on matrices of order at most \(p\) all use polynomially many graph or exact rational-arithmetic operations.  This proves the runtime assertion.

%%% FIELD proof-finite-dictionary
If \(\mathsf S(C)\ne\mathsf S(D)\), then \(\texttt{thm:bounded-binary-star-separation}\) returns a source \(E\in\{C,D\}\) and a separating covariance activating one or two blocks at one target.  By \(\texttt{def:binary-star-library}\), it lies in \(\mathcal W_2(E)\) and not in the other model.  Hence
\[
 \mathcal W_2(C)\not\subseteq\mathcal M(D)
 \quad\text{or}\quad
 \mathcal W_2(D)\not\subseteq\mathcal M(C).
\]
For \(b_j=|\Pi^{E_0}_j|\), target \(j\) offers \(b_j+\binom{b_j}{2}\) choices.  Duplicate covariances only reduce the cardinality.  Since \(b_j\le|\operatorname{pa}_{G_{E_0}}(j)|\le p-1\),
\[
 |\mathcal W_2(E_0)|
 \le\sum_{j\in V}\frac{b_j(b_j+1)}2
 \le p\frac{(p-1)p}{2}.
\]
Empty parent sets contribute zero.

%%% FIELD statement-chickering
Let \(G_0\) and \(G_1\) be finite directed acyclic graphs on the same labelled vertex set.  They have the same skeleton and the same unshielded colliders if and only if they are Markov equivalent.  If they are equivalent and \(m\) adjacencies have opposite orientations, then there is a sequence of exactly \(m\) distinct covered-edge reversals from \(G_0\) to \(G_1\); every intermediate graph is acyclic and equivalent, and each reversal corrects one original orientation disagreement.

%%% FIELD proof-complete-equivalence
If \(\mathcal M(C)=\mathcal M(D)\) but \(\mathsf S(C)\ne\mathsf S(D)\), the bounded separation theorem supplies \(\Sigma^{\star}\) in exactly one of the two models, a contradiction.  Thus model equality implies signature equality.

Assume \(\mathsf S(C)=\mathsf S(D)\).  By \(\texttt{def:static-signature}\), \(G,H\) have the same skeleton and colliders.  The cited \(\texttt{lem:chickering-covered-disagreement-sequence}\) gives
\[
 G=G_0,G_1,\ldots,G_m=H,\qquad m=\delta(C,D),
\]
where each step reverses a covered edge that disagreed between \(G,H\), and no such edge is reversed twice.  Write \(\mathcal N_j=\mathcal N_j(C)=\mathcal N_j(D)\).  Every arrow \(a\to j\) with \(a\in B\in\mathcal N_j\) occurs at both endpoints, so it is never reversed.  Lift each \(G_r\) by
\[
 \Pi^{(r)}_j=\mathcal N_j\cup
 \{\{a\}:a\in\operatorname{pa}_{G_r}(j)\setminus\bigcup\mathcal N_j\}.
\]
This is a valid partition, and at \(r=0,m\) it equals the partition of \(C,D\), respectively.

Suppose the next covered disagreement is \(i\to j\).  Equality of \(F(C),F(D)\) forces \(\mathcal N_i=\mathcal N_j\), since otherwise the disagreeing orientation would be recorded in \(F\).  Coveredness says
\[
 \operatorname{pa}_{G_r}(j)=\operatorname{pa}_{G_r}(i)\cup\{i\}.
\]
The parent \(i\) cannot belong to a block in \(\mathcal N_j=\mathcal N_i\), for that would make a parent block of \(i\) contain \(i\).  Hence \(\{i\}\in\Pi^{(r)}_j\).  Removing the common nonsingleton blocks from the covered parent-set equality leaves identical singleton blocks, so
\[
 \Pi^{(r)}_j\setminus\{\{i\}\}=\Pi^{(r)}_i.
\]
The move is legal and its prescribed transport produces \(\Pi^{(r+1)}\).  Induction proves \(D\in\mathcal R(C)\).

It remains to prove that one legal reversal preserves the full image.  Put \(S=\operatorname{pa}(i)\), so \(\operatorname{pa}(j)=S\cup\{i\}\), and write
\[
 X_i=a^{\mathsf T}X_S+\varepsilon_i,\qquad
 X_j=bX_i+d^{\mathsf T}X_S+\varepsilon_j,\qquad
 \operatorname{Var}(\varepsilon_i)=u>0,\quad
 \operatorname{Var}(\varepsilon_j)=v>0.
\]
Legality makes \(a,d\) block-constant on the same partition.  Define
\[
 t=v+b^2u>0,\quad q=\frac{bu}{t},\quad
 \eta_j=b\varepsilon_i+\varepsilon_j,\quad
 \eta_i=\varepsilon_i-q\eta_j.
\]
Substitution gives
\[
 X_j=(d+ba)^{\mathsf T}X_S+\eta_j,\qquad
 X_i=qX_j+\{a-q(d+ba)\}^{\mathsf T}X_S+\eta_i.
\]
Moreover,
\[
 \operatorname{Var}(\eta_j)=t,\quad
 \operatorname{Cov}(\eta_i,\eta_j)=bu-qt=0,\quad
 \operatorname{Var}(\eta_i)=u-\frac{b^2u^2}{t}=\frac{uv}{t}>0.
\]
The transformed innovations are jointly Gaussian, hence their zero covariance gives independence; they remain independent of all other innovations.  The new vectors are linear combinations of block-constant vectors, and \(q\) belongs to the new singleton block.  This is a valid parameter of the reversed state with the same law.  The reverse move is legal as well, proving equality of images in both directions, including \(b=0\).  Iteration proves \(D\in\mathcal R(C)\Rightarrow\mathcal M(C)=\mathcal M(D)\).

The constructed path has length \(\delta(C,D)\), so \(d_{\mathrm{rev}}(C,D)\le\delta(C,D)\).  Every path must reverse each initially disagreeing adjacency an odd number of times, giving the opposite inequality.  Hence equality holds, and the constructed shortest path reverses each disagreement once.

%%% FIELD statement-complete-equivalence
For all \(C,D\in\mathfrak S_p\), the following are equivalent: \(\mathcal M(C)=\mathcal M(D)\); \(\mathsf S(C)=\mathsf S(D)\); and \(D\in\mathcal R(C)\). Every legal reversal preserves the full covariance image, including zero coefficients. If these conditions hold, then \(d_{\mathrm{rev}}(C,D)=\delta(C,D)\), and a shortest path reverses every initially disagreeing arrow exactly once. Thus the signature is an exact recovery functional for the covariance-model component, not merely a sufficient invariant.

%%% FIELD proof-enumeration
If \(B\in\mathcal N_j(C)\) and \(i\in B\), then \(\mathcal N_i(C)\ne\mathcal N_j(C)\); equality would make \(B\) a parent block of \(i\) containing \(i\).  Thus every arrow supporting a nonsingleton block belongs to \(F(C)\).  Every extension \(K\) in \(\texttt{def:component-enumerator}\) retains these arrows, so
\[
 \Pi^K_j=\mathcal N_j(C)\cup
 \{\{a\}:a\in\operatorname{pa}_K(j)\setminus\bigcup\mathcal N_j(C)\}
\]
is a valid partition.  It is the unique partition with nonsingleton family \(\mathcal N_j(C)\), since every residual parent must be singleton.

Each emitted \(K\) has the skeleton and colliders of \(G\), retains \(F(C)\), and has the same endpoint families.  Therefore its lifted state \(E_K\) has \(\mathsf S(E_K)=\mathsf S(C)\).  Conversely, a state with this signature is an acyclic extension with the required colliders and background arrows, and uniqueness of the lift makes it an output.  The complete-equivalence theorem now gives
\[
 \operatorname{Enum}(C)=\{D:\mathsf S(D)=\mathsf S(C)\}=\mathcal R(C).
\]
The unique lift makes the output duplicate-free.  Choose the lexicographically first topological order of every emitted DAG.  Two orientations of a fixed skeleton cannot share a topological order, because every edge points from its earlier to its later endpoint.  This injects the outputs into the \(p!\) total orders, proving termination after at most \(p!\) outputs.  Finally, the complete-equivalence theorem supplies a shortest legal path of length \(\delta(C,D)\) to each output.

%%% FIELD proposed-ordinary-proper
If every block of \(C\) is a singleton, then \(\mathcal N_j(C)=\varnothing\) for every \(j\in V\), and graph projection is a bijection from \(\mathcal R(C)\) onto the ordinary Markov-equivalence class of \(G\); the inverse sends each equivalent DAG to its unique all-singleton lift.  If every nonempty block of \(C\) is nonsingleton, then \(\mathcal R(C)=\{C\}\).  The first conclusion is the singleton lift of the Chickering1995 covered-reversal component, and the second is the target-local proper-blocked structural-identifiability conclusion covered by BoegeKubjasMisraSolus2026.

%%% FIELD proposed-ordinary-proper-postimage
If every block of \(C\) is a singleton, then \(\mathcal N_j(C)=\varnothing\) for every \(j\in V\), and graph projection is a bijection from \(\mathcal R(C)\) onto the ordinary Markov-equivalence class of \(G\); the inverse sends each equivalent DAG to its unique all-singleton lift.  If every nonempty block of \(C\) is nonsingleton, then \(\mathcal R(C)=\{C\}\).  The first conclusion is the singleton lift of the Chickering1995 covered-reversal component, and the second is the target-local proper-blocked structural-identifiability conclusion covered by BoegeKubjasMisraSolus2026.

%%% FIELD proof-ordinary-proper
If every block of \(C\) is singleton, then \(\texttt{def:nonsingleton-family}\) gives \(\mathcal N_j(C)=\varnothing\) for all \(j\), and \(\texttt{def:family-mismatch-edges}\) gives \(F(C)=\varnothing\).  Every DAG \(H\) ordinarily Markov equivalent to \(G\) has a unique all-singleton lift \(D_H\); it has the same skeleton and colliders and empty nonsingleton families and \(F\)-set.  Thus \(\mathsf S(D_H)=\mathsf S(C)\), so the complete-equivalence theorem places \(D_H\) in \(\mathcal R(C)\).  Conversely, every member of \(\mathcal R(C)\) has that signature, hence is an all-singleton lift of a DAG ordinarily equivalent to \(G\).  This proves the claimed bijection, and the cited Chickering lemma identifies its graph image with the covered-reversal component.

If every nonempty block of \(C\) is nonsingleton, a legal move out of \(C\) would require a singleton block, contrary to the hypothesis.  A legal move cannot enter \(C\), either, because its endpoint partition contains the new singleton specified by \(\texttt{def:legal-reversal}\).  Thus \(C\) is isolated and \(\mathcal R(C)=\{C\}\), agreeing with the target-local specialization of \(\texttt{lem:boege-proper-blocked-identifiability}\).

%%% FIELD proof-bic
Fix \(E_0\).  The columns of \(A_{E_0j}\) are nonzero indicators of disjoint blocks, hence linearly independent.  Thus for \(c\ne0\) and \(\Sigma\in\mathbb S^p_{++}\),
\[
 c^{\mathsf T}\Gamma_{E_0j}(\Sigma)c
 =(A_{E_0j}c)^{\mathsf T}\Sigma_{P,P}(A_{E_0j}c)>0.
\]
Every inverse in \(\texttt{def:off-model-block-ls}\) exists.  Its constrained residual variance is positive because
\[
 \min_b\operatorname{Var}\{X_j-b^{\mathsf T}A_{E_0j}^{\mathsf T}X_P\}
 \ge\Sigma_{jj}-\Sigma_{jP}\Sigma_{P,P}^{-1}\Sigma_{Pj}>0.
\]

On \(\mathcal M(E_0)\), innovation independence makes parent regression recover every structural row and innovation variance.  These quantities form a smooth rational inverse to \(\phi_{E_0}\).  Therefore \(\mathcal M(E_0)\) is a connected smooth embedded manifold of dimension \(d_{E_0}=p+\sum_j|\Pi^{E_0}_j|\).  In a topological order, membership is equivalent to requiring, in each regression on all predecessors, zero coefficients for nonparents and equal coefficients within each block.  Necessity follows from the triangular SEM; conversely, the ordered Gaussian conditional densities yield such an SEM with positive conditional variances.  Clearing the positive principal-minor denominators makes these finitely many polynomial equations.  Consequently \(\mathcal M(E_0)\) is relatively closed in \(\mathbb S^p_{++}\).

If \(d_D<d_C\), the \(d_C\)-dimensional manifold \(\mathcal M(C)\) cannot lie in \(\mathcal M(D)\).  If \(d_D=d_C\) and containment held, equal-dimensional charts would make \(\mathcal M(C)\) open in \(\mathcal M(D)\); relative closedness makes it closed, and connectedness of \(\mathcal M(D)\) would force equality.  Thus every nonequivalent \(D\) with \(d_D\le d_C\) has a defining polynomial whose pullback through \(\phi_C\) is nonzero.  A nonconstant rational effect coordinate has some nonzero partial derivative; clearing positive denominators makes its zero-gradient locus proper algebraic.  Finiteness of the state and query sets, and multiplication of one nonzero polynomial for each locus, show that \(\mathcal E_C\) is a proper algebraic subset of \(\Theta_C\).

Let \(\theta_C\notin\mathcal E_C\) and \(\Sigma_0=\phi_C(\theta_C)\).  For \(K\in\mathbb S^p_{++}\), put
\[
 q_T(K)=\log\det K+\operatorname{tr}(K^{-1}T).
\]
Up to a state-independent constant, \(-2n^{-1}\ell_n(K)=q_{S_n}(K)\).  Gaussian fourth moments give \(\mathbb E\|S_n-\Sigma_0\|_F^2=O(n^{-1})\); hence
\[
 S_n\to\Sigma_0,\qquad \|S_n-\Sigma_0\|_F=O_{\mathbb P}(n^{-1/2}).
\]
For \(n\ge p\), the Gaussian data matrix has full row rank almost surely, because every rank-deficient matrix lies in the zero set of its nonzero \(p\)-minor polynomials.  Thus the singular fallback is eventually irrelevant.

For nonequivalent \(D\) with \(d_D\le d_C\), exclusion from \(\mathcal E_C\) gives \(\Sigma_0\notin\mathcal M(D)\).  If \(\lambda_r\) are the eigenvalues of \(K^{-1}\Sigma_0\), then
\[
 q_{\Sigma_0}(K)-q_{\Sigma_0}(\Sigma_0)
 =\sum_{r=1}^p(\lambda_r-\log\lambda_r-1)\ge0,
\]
with equality only at \(K=\Sigma_0\).  The summand diverges at zero and infinity, so sublevel sets are compact.  Relative closedness yields
\[
 \Delta_D=\inf_{K\in\mathcal M(D)}
 \{q_{\Sigma_0}(K)-q_{\Sigma_0}(\Sigma_0)\}>0.
\]
Uniform continuity on a compact sublevel set and \(S_n\to\Sigma_0\) make the empirical gap at least \(\Delta_D/2\) with probability tending to one.  Since \(\Sigma_0\in\mathcal M(C)\),
\[
 \operatorname{BIC}_n(D)-\operatorname{BIC}_n(C)
 \ge n\Delta_D/2-(d_C-d_D)\log n\to+\infty.
\]

If \(d_D>d_C\), then
\[
 \sup_{\mathcal M(D)}\ell_n-\sup_{\mathcal M(C)}\ell_n
 \le\sup_{\mathbb S^p_{++}}\ell_n-\ell_n(\Sigma_0).
\]
The saturated maximizer is \(S_n\).  A second-order expansion at \(\Sigma_0\) and the preceding \(n^{-1/2}\) bound make the right side \(O_{\mathbb P}(1)\).  Therefore
\[
 \operatorname{BIC}_n(D)-\operatorname{BIC}_n(C)
 \ge-O_{\mathbb P}(1)+(d_D-d_C)\log n\to+\infty.
\]
If \(\mathcal M(D)=\mathcal M(C)\), the complete-equivalence theorem connects the states by legal reversals.  Each move removes one singleton block and adds one, so \(d_D=d_C\), and the BIC values tie exactly.  Since \(\mathfrak S_p\) is finite for fixed \(p\), a union bound over nonequivalent states proves
\[
 \Pr\{\mathcal R(\widehat C_n)=\mathcal R(C)\}\longrightarrow1.
\]
Only the global criterion was used.

%%% FIELD proof-effect-inference
Fix \(E\in\mathcal R(C)\).  Complete equivalence gives \(\Sigma_0\in\mathcal M(E)\).  If the structural parent vector at \(j\) is \(A_{Ej}\beta_{Ej}\), innovation independence gives
\[
 \Sigma_{P,j}=\Sigma_{P,P}A_{Ej}\beta_{Ej},\qquad
 h_{Ej}(\Sigma_0)=\Gamma_{Ej}(\Sigma_0)\beta_{Ej}.
\]
The positive-definite Gram matrix is invertible, so \(L_E(\Sigma_0)=B_E\).  Acyclicity makes \(B_E\) nilpotent.  For \(a\ne y\), differentiating the intervened SEM after replacing the equation for \(X_a\) by \(X_a=x\) gives
\[
 \tau^E_{a\to y}=\sum_{k=1}^{p-1}(B_E^k)_{ya}
 =((I_p-B_E)^{-1})_{ya}=g_E(\Sigma_0)_{(a,y)}.
\]
Thus equality of the causal effect and observed-law functional is proved from the structural equations; it is not definitional.  Since the observed covariance identifies only \(\mathcal R(C)\), the member-labelled effects form the set in \(\texttt{def:component-effect-set}\).

For arbitrary positive-definite \(\Sigma\), every block Gram matrix remains positive definite, its inverse is rational and smooth, and
\[
 (I_p-L_E(\Sigma))^{-1}=\sum_{k=0}^{p-1}L_E(\Sigma)^k.
\]
Thus every \(g_E\) is smooth near \(\Sigma_0\).

We derive the covariance limit used by the intervals.  For symmetric \(T\), direct Gaussian integration gives
\[
 \mathbb E\exp\{isX^{\mathsf T}TX\}
 =\det(I_p-2is\Sigma_0T)^{-1/2}.
\]
Independence and \(\log\det(I-U)=-\operatorname{tr}U-\tfrac12\operatorname{tr}(U^2)+O(\|U\|^3)\) yield
\[
 \mathbb E\exp\!\left[i\operatorname{tr}\{T\sqrt n(S_n-\Sigma_0)\}\right]
 \longrightarrow\exp\{-\operatorname{tr}((\Sigma_0T)^2)\}.
\]
This is the characteristic function of a centered Gaussian symmetric matrix \(Z\) with
\[
 \operatorname{Cov}(Z_{ij},Z_{k\ell})
 =\Sigma_{0,ik}\Sigma_{0,j\ell}+\Sigma_{0,i\ell}\Sigma_{0,jk}.
\]
For nonzero symmetric \(T\), the variance is
\[
 2\operatorname{tr}\{(\Sigma_0^{1/2}T\Sigma_0^{1/2})^2\}>0.
\]

Let \(g\) be a nonconstant labelled coordinate.  Since \(\theta_C\notin\mathcal E_C\), \(Dg(\Sigma_0)\ne0\).  Taylor expansion gives
\[
 \sqrt n\{g(S_n)-g(\Sigma_0)\}
 =Dg(\Sigma_0)[\sqrt n(S_n-\Sigma_0)]+o_{\mathbb P}(1)
 \Longrightarrow N(0,\sigma_g^2),
\]
with \(\sigma_g^2>0\).  The plug-in variance is consistent by continuity, so its studentized interval has limiting coverage \(1-\alpha/K\).  A structurally zero coordinate has no directed path, is identically zero by the finite path expansion, and its singleton interval covers exactly.  There are no other constant coordinates: zero block coefficients give effect zero, while positive coefficients on a directed path give a positive effect.

A union bound over the \(K\) nonconstant labels gives simultaneous limiting coverage at least \(1-\alpha\), without a nonsingular joint effect covariance.  Generic BIC recovery makes the probability of selecting the wrong component \(o(1)\), proving the stated liminf after intersecting the two events.  Labels are retained, so duplicate values are harmless.

On the correct-component event, continuity and finiteness give
\[
 \max_{E\in\mathcal R(C)}\|g_E(S_n)-g_E(\Sigma_0)\|\longrightarrow0
\]
in probability.  The Hausdorff distance of the corresponding unlabelled finite sets is bounded by this maximum.  Adding the \(o(1)\) selection-error probability proves the asserted Hausdorff convergence.

%%% FIELD statement-effect-inference
Under \(\texttt{thm:generic-bic-component-recovery}\), fix finite \(\mathcal Q\) and \(\alpha\in(0,1)\). Then \(\mathcal T_{\widehat C_n}(S_n)\) converges in probability to \(\mathcal T_C(\Sigma_0)\) in Hausdorff distance. Retaining state labels before removing coincident numerical values, the confidence collection satisfies
\[
\liminf_{n\to\infty}\Pr\left\{\mathcal R(\widehat C_n)=\mathcal R(C)\text{ and every }g_E(\Sigma_0),\ E\in\mathcal R(C),\text{ belongs to its labelled interval in }\mathfrak I_{n,\alpha}(\widehat C_n)\right\}\geq1-\alpha.
\]
The conclusion permits duplicate effect values and singular joint effect covariance.

%%% FIELD partial-pareto
Within the frozen binary-star construction, \(\texttt{thm:bounded-binary-star-separation}\) supplies the upper point \((2,\lceil\log_2p\rceil)\).  For every \(p\ge2\), any source-SEM separator must activate at least one block: no activation with unit innovations returns \(I_p\), which belongs to every state model by setting all structural coefficients to zero, while unequal signatures exist between the empty graph and a one-edge graph.  Hence the first coordinate is at least \(1\), and height is nonnegative.  The frozen core does not decide whether one block can be universal or whether height below \(\lceil\log_2p\rceil\) is possible.
